% Replacement for Section 4. Uses the manuscript's existing math/theorem packages.
% The accompanying appendices preserve the detailed derivations and proofs.
\section{GOTHAM for Wiener systems with AR output noise}
\label{sec:GOTHAM_for_WS_AR_noise}

This section specializes the operator formulation of
Section~\ref{subsec:OHAM_reformulation} to the Wiener--AR model of
Section~\ref{sec:problem_formulation}. Only the measured input $u$, the output
record $C$, and reconstructed quantities enter the identification criteria;
the true innovations are not available to the estimator. At each outer
iteration, the current estimate initializes the homotopy and the first $M$
deformation terms are summed according to
\eqref{eq:NOHAM_iterative_update}. The explicit rules below cover both
parametrizations, avoiding separate derivations for each configuration.

\subsection{Identification criteria}
\label{subsec:GOTHAM_PL_parametrization}
\label{subsec:GOTHAM_FSM_parametrization}

For the pseudo-linearized (PL) parametrization, the reconstructed regressor
$\hat\Phi^{(k)}$ is fixed throughout the current homotopy construction:
\begin{equation}
\label{eq:WS_loss_k}
    J_{\mathrm{PL}}^{(k)}(\theta)
    =\frac12\|C-\hat\Phi^{(k)}\theta\|_2^2.
\end{equation}
Its gradient is affine and its Hessian is
$G^{(k)}=\hat\Phi^{(k)\top}\hat\Phi^{(k)}$. The latent signals and regressor
are reconstructed again before the next outer iteration.

For the direct nonlinear FSM parametrization, let
$\beta_{\mathrm{FSM}}(t;\theta)$ denote the signal-path reconstruction
defined through the frequency-sampled impulse-response coefficients
in~\eqref{eq:fsm_impulse_coefficients}. The criterion is
\begin{align}
\label{eq:FSM_innovation_residual_filter}
    r_{\nu,\mathrm{FSM}}(t;\theta)
    &=D(z^{-1};\theta)[c(t)-\beta_{\mathrm{FSM}}(t;\theta)],\\
\label{eq:FSM_innovation_energy_loss}
    J_{\nu,\mathrm{FSM}}(\theta)
    &=\frac12\sum_{t=t_0}^{\tau}r_{\nu,\mathrm{FSM}}^2(t;\theta).
\end{align}
Both the signal path and the candidate AR polynomial are differentiated
with respect to $\theta$; $t_0$ accounts for the chosen initialization or
transient exclusion. This innovation-energy criterion avoids the algebraic
cancellation obtained by reconstructing an innovation and passing it back
through its matching noise filter; see
Remark~\ref{rem:fsm_innovation_energy} and
Appendix~\ref{app:fsm_innovation_energy}.

Within one outer iteration, suppress the superscript $(k)$ and write
$J$ for the selected criterion, $\hat\theta$ for the current estimate, and
\begin{equation}
\label{eq:GOTHAM_local_quantities}
    g=\nabla J(\hat\theta),\qquad
    Q=\nabla^2J(\hat\theta),\qquad
    D_Q=\operatorname{diag}(\operatorname{diag}Q).
\end{equation}
Thus $Q=G^{(k)}$ for PL, whereas FSM uses the full nonlinear Hessian.
In all configurations below, the HAM residual-space operator is $H=I$
and $\hbar=-\eta$, with $\eta>0$.

\subsection{Gradient-root updates and classical special cases}
\label{subsec:GOTHAM_PL_gradient_root}
\label{subsec:GOTHAM_update_summary}

For $\mathcal N_G=\nabla J$, let $P=\mathcal L^\dagger$ be the fixed
linear solution operator during the current homotopy construction.
The first two corrections are
\begin{equation}
\label{eq:GOTHAM_common_gradient_corrections}
    \theta_1=-\eta Pg,\qquad
    \theta_2=\theta_1-\eta PQ\theta_1.
\end{equation}
Consequently, the order-one and order-two increments are
$\Delta_1=\theta_1$ and $\Delta_2=\theta_1+\theta_2$, respectively, with
$\hat\theta^{+}=\hat\theta+\Delta_M$.
Table~\ref{tab:GOTHAM_update_rules} gives their explicit forms. These
expressions hold for either criterion; higher homotopy orders generally
require higher derivatives for FSM.

\begin{table}[t]
\centering
\small
\setlength{\tabcolsep}{7pt}
\renewcommand{\arraystretch}{1.45}
\caption{Gradient-root GOTHAM increments for either parametrization.
The inverse symbols in this table denote Moore--Penrose inverses.
For regularized operators, use the corresponding $P$ directly
in~\eqref{eq:GOTHAM_common_gradient_corrections}.}
\label{tab:GOTHAM_update_rules}
\begin{tabular}{@{}ccl@{}}
\toprule
$M$ & $\mathcal L$ & Increment $\Delta_M$\\
\midrule
1 & $I$   & $-\eta g$\\
1 & $Q$   & $-\eta Q^\dagger g$\\
1 & $D_Q$ & $-\eta D_Q^\dagger g$\\
\midrule
2 & $I$   & $-(2\eta I-\eta^2Q)g$\\
2 & $Q$   & $-(2\eta-\eta^2)Q^\dagger g$\\
2 & $D_Q$ & $-D_Q^\dagger(2\eta g-\eta^2QD_Q^\dagger g)$\\
\bottomrule
\end{tabular}
\end{table}

For PL, the first-order identity choice recovers the WS-GGI gradient
update of~\cite{Lv2026}, with the step normalization absorbed into $\eta$.
The first-order full-Hessian choice is the fixed-regressor
Gauss--Newton update, denoted WS-GNI here. If $G^{(k)}$ is nonsingular
and $\eta=1$, it gives the batch least-squares solution
$(G^{(k)})^{-1}\hat\Phi^{(k)\top}C$. For FSM, the corresponding choices
are gradient and Newton-type updates. The second-order full-Hessian
row is a scaled first-order step for both parametrizations when the
exact Moore--Penrose inverse is used. Proofs, the rank-deficient case,
and the effect of regularization are given in
Appendix~\ref{app:gotham_pl_updates}.

\begin{remark}[Choice of auxiliary operator]
The identity, diagonal, and full-Hessian choices trade computational cost
against curvature information and conditioning. Regularization changes
the update and can invalidate simplified equivalences in
Table~\ref{tab:GOTHAM_update_rules}; an arbitrary Hessian shift also need
not produce a descent direction when the FSM Hessian is indefinite.
These choices and the convergence-control parameter must therefore be
specified together.
\end{remark}

\subsection{Loss-root updates and the Polyak connection}
\label{subsec:GOTHAM_PL_loss_root}

For $\mathcal N_L=J-\varepsilon_0$, take
$\mathcal Lh=g^\top h$. Its Moore--Penrose or Tikhonov solution
operator is $s\mapsto v_\lambda s$, where
\begin{equation}
\label{eq:GOTHAM_common_loss_inverse}
    v_\lambda=\frac{g}{\|g\|_2^2+\lambda}.
\end{equation}
Here $\lambda=0$ requires $g\ne0$, and $\lambda>0$ corresponds to
a penalty $\lambda\|h\|_2^2$ in the linear subproblem. For either PL or
FSM, the first three corrections are
\begin{equation}
\label{eq:GOTHAM_common_loss_corrections}
\begin{aligned}
    \theta_1&=-\eta v_\lambda[J(\hat\theta)-\varepsilon_0],\\
    \theta_2&=\theta_1-\eta v_\lambda g^\top\theta_1,\\
    \theta_3&=\theta_2-\eta v_\lambda
       \left(g^\top\theta_2+\tfrac12\theta_1^\top Q\theta_1\right).
\end{aligned}
\end{equation}
Summing the first $M$ terms gives the order-$M$ update, for $M=1,2,3$.
With $M=1$, $\lambda=0$, and $\eta=1$, the step size is
$[J(\hat\theta)-\varepsilon_0]/\|g\|_2^2$. It is the classical
Polyak rule~\cite{Polyak1969} when $\varepsilon_0=J_\star$, the optimal
loss value. A prescribed or estimated target gives a target-level variant;
the same convergence guarantees do not follow merely from this algebraic
similarity. The derivations and precise target conditions are given in
Appendix~\ref{app:gotham_loss_updates}.

\begin{remark}[Meaning of the loss target]
An exact loss root requires an attainable level
$\varepsilon_0\in J(\Theta)$; nonnegativity alone is insufficient.
For $\varepsilon_0>J_\star$, attaining the target need not imply
$g=0$ or a unique parameter estimate. Since the residual is scalar,
its deformation equations are underdetermined when $n>1$, and the
chosen inverse determines how the parameter correction is selected.
\end{remark}

\begin{remark}[Implementation and method names]
For PL, the quadratic structure gives explicit residual coefficients.
For FSM, automatic differentiation supplies the required derivatives
without expanding the nested frequency-response and polynomial terms;
see Appendix~\ref{app:gotham_fsm_updates}. In
Section~\ref{sec:results}, GHAM and LHAM denote gradient-root and
loss-root instances, the integer specifies $M$, and $I$, $H$, $dH$,
and TIK identify the auxiliary operator or inversion strategy.
The method-name suffix $H$ denotes a Hessian, whereas the HAM
residual-space operator remains $H=I$ throughout this section.
\end{remark}
